\documentclass{article}
\usepackage[T1]{fontenc}
\usepackage[utf8]{inputenc}
\usepackage{parskip}

\setlength{\parskip}{8pt}
% Linespread command allows you to change line spacing for the entire document
\linespread{1.18}

% Tweak page margins
\addtolength{\oddsidemargin}{-.875in}
\addtolength{\evensidemargin}{-.875in}
\addtolength{\textwidth}{1.75in}

\addtolength{\topmargin}{-.875in}
\addtolength{\textheight}{1.75in}

\usepackage{natbib}
\usepackage{hyperref}
\usepackage{xcolor}
\usepackage{xspace}
\usepackage{fancyhdr}
\hypersetup{
    colorlinks,
    linkcolor={red!50!black},
    citecolor={blue!50!black},
    urlcolor={blue!80!black}
}

\newcommand{\HRule}{\rule{\linewidth}{0.5mm}}
\newcommand{\Hrule}{\rule{\linewidth}{0.3mm}}

% Project specific macros
\newcommand{\graphite}{GRAPHITE\xspace}
\newcommand{\wave}{WAVE\xspace}

% School specific macros
\newcommand{\schoolShort}{NU\xspace}
\newcommand{\school}{Northwestern University\xspace}
\newcommand{\schoolLong}{Northwestern University\xspace}

\newcommand{\profOne}{Prof. Ryan Huang\xspace}
\newcommand{\profTwo}{Prof. Xiaonan Huang\xspace}
\newcommand{\profThree}{Prof. Yutong Ban\xspace}

% Creates header for each page
\usepackage{fancyhdr}
\pagestyle{fancy}
\fancyhf{}
\fancyhead[LE,RO]{\header\hskip\linepagesep\vfootline\thepage}
\newskip\linepagesep \linepagesep 5pt\relax
\def\vfootline{%
    \begingroup
    	\rule[-10pt]{0.75pt}{25pt}
    \endgroup
}
\def\header{%
	\begin{minipage}[]{120pt}
		\hfill Yuchen You
    	\par \hfill 				% Formatting boilerplate
    	MSCS, Fall 2026 			% Area, Program, Cycle, Year
    \end{minipage}
}
\fancyhead[RE,LO]{Academic Statement | \schoolLong}
\renewcommand\headrulewidth{0pt}

\begin{document}

\paragraph{Research Focus and Motivation.}
I am applying to the MS in Computer Science (MSCS) program at \school{} to study \textbf{reliability for distributed systems and ML serving} by building \textbf{telemetry-to-action} mechanisms: given noisy, partial signals (metrics, logs, traces, and alerts), my goal is to infer SLO-relevant system state and trigger mitigations that are \textbf{safe by construction} (e.g., stability and “never-worsen” constraints). My recent work on ZooKeeper overload mitigation and launch-bound GPU runtimes convinced me that today's stacks still expose telemetry and knobs without a principled interface for \emph{when} and \emph{how} to act---leaving reliability dependent on manual interpretation and ad hoc playbooks.

\paragraph{Professional Plans and the Role of a Master.}
In the long term, I aim to become a \textbf{systems engineer and researcher} who designs and evaluates \textbf{reliable infrastructure for ML and cloud workloads}, where correctness and latency behavior matter under real failures and multi-tenancy. The MSCS program is essential training for me to turn this direction into durable contributions: to formulate sharp research questions, develop abstractions that couple observability with control under explicit safety constraints, and validate them through rigorous experimentation and realistic fault models---so reliability improves systematically rather than by operator heroics.

\paragraph{Reliability Lessons from Embedded Systems.}
Before focusing on large-scale systems, advised by \textbf{Prof.~Xiaonan~Huang}, I worked on embedded and robotics; our modular robotic arm project was selected for a \textbf{Best Poster Award} at an \textbf{ICRA workshop}. That experience made reliability concrete: failures often came from noisy or delayed signals and weak fault handling rather than from control laws alone. Simple monitors, sanity checks, and degraded-but-safe modes were often the difference between a system that merely worked in ideal conditions and one that remained usable under stress. This perspective now shapes how I approach distributed and ML systems: prioritize observability, fault containment, and graceful degradation.

\paragraph{Agentic Distributed System Operations.}
At \textbf{UMich's OrderLab} with \textbf{Prof.~Ryan Huang}, I study safe automated mitigation for overload and gray failures in ZooKeeper under partial observability, often relying on coarse anomaly signals.
Our question is whether an agentic controller can match a tuned rule-based baseline on \emph{aggregate client throughput} and \emph{avg/p95 latency} without inducing oscillation or excessive throttling.
To support controlled evaluation, I built an end-to-end \textbf{ZooKeeper testbed} with Prometheus telemetry and HAProxy-based shaping, together with a \textbf{typed mitigation-action library} whose execution is guarded by explicit safety checks, and a fault-injection/benchmarking harness (\texttt{zkbench}) that supports weighted injections and load-peak scenarios with repeated trials to control variance.
As a baseline, we implemented a \textbf{rule-based controller} that triggers actions via \textbf{sliding-window trend checks} and executes parameterized mitigations from a YAML playbook mapping failure modes to actions. While it prevents collapse, it often over- or under-reacts to workload shifts, leaving extended tail-latency windows.
Building on this substrate, I developed an \textbf{agentic layer} where an LLM proposes concrete mitigation plans from the same constrained action set based on recent metrics and finite-state summaries; early failures with a richer API (e.g., conflicting actions and over-throttling) pushed me to tighten the observation/action interface and formalize guardrails.
We therefore treat the LLM as a \textbf{proposal generator} and gate execution with explicit safety bounds and guardrail checks based on recent telemetry.
With these guardrails, the agent consistently shortened SLO-violation windows under gray failures and skewed workloads, reinforcing my view that the key to reliable closed-loop control is rarely ``smarter decisions'' alone, but \textbf{well-chosen state abstractions and guardrails} that keep decisions stable under noisy signals. Looking ahead, I plan to stress the system under more complex failure compositions and refine mitigation selection while preserving these safety guarantees.

\paragraph{CUDA Proxy Player: Runtime Support for Launch-Bound Inference.}
In an Advanced Operating Systems course project, we asked when conditional GPU inference (e.g., Mixture-of-Experts serving) truly becomes \emph{launch-bound}, with host-side orchestration dominating end-to-end latency. Based on the hypothesis that \emph{shape-stable} regions can be amortized while \emph{unstable} glue should be isolated, I co-designed CUDA Proxy Player: a multi-path runtime that routes requests by size and shape to eager, CUDA Graph replay, or persistent-kernel execution, and stitches these flows together via explicit capture/replay and scheduling interfaces. I built the benchmarking framework---microbenchmarks and end-to-end experiments spanning expert sizes, batch shapes, and traffic mixes---to quantify where each path wins and why. In the most launch-bound regimes, size-aware routing enabled the CUDA Graph path to deliver substantial gains over a naïve eager baseline; however, our breakdown also showed why a hybrid graph+worker design often fails to translate into end-to-end wins: coordination overheads (worker scheduling and phase barriers) can dominate, and keeping workers/buffers resident reduces VRAM headroom and can worsen tail latency under multi-tenant workloads. Designing these experiments taught me to separate real speedups from coordination artifacts, and reinforced my view that ML serving runtimes should expose telemetry and control hooks for VRAM usage, admission control, and path selection---so operators can tune policies to deployment constraints rather than chasing a single ``fastest'' path.

\paragraph{Fit with \schoolShort's Systems Curriculum and Course-Based MSCS Plan.}
\school is an ideal place for me to deepen systems expertise through a rigorous, course-based MSCS curriculum that still emphasizes \emph{building} and \emph{evaluation}.
\textbf{COMP\_SCI 310 (Scalable Software Architectures)} will help me translate my ZooKeeper reliability work into end-to-end service design skills by studying real-world high-scale systems and implementing cloud-facing components, while \textbf{COMP\_SCI 468 (Programming Massively Parallel Processors with CUDA)} will strengthen the GPU performance-engineering toolkit behind my CUDA runtime work through hands-on CUDA programming and an open-ended final project.
Complementing these, I also plan to strengthen my core systems foundation through advanced coursework in operating systems, distributed systems, and networking, so my project work is grounded in first principles rather than ad hoc tuning.
Within this course-based track, I will pursue a coherent set of systems and ML-infrastructure projects that connect \emph{observability} to \emph{action}: designing constrained, safety-checked mitigation interfaces for distributed services and analyzing how runtime choices (e.g., CUDA graphs vs.\ persistent execution) shape tail latency under contention.
By combining these implementation-heavy experiences with disciplined measurement and reproducible evaluation, \school's MSCS program will prepare me to build reliable, high-performance infrastructure for modern distributed and ML serving systems.

\end{document}

% That's All Folks.

% Best of luck, you got this! :)